\subsection{Context and Motivations}
% Descrivere il contesto scientifico e tecnologico in cui si colloca l'attività di tesi e quali sono le motivazioni per lo sviluppo del progetto di tesi. (circa 2 pagine)
The constant and rapid evolution in the field of Artificial Intelligence has result in a day-to-day use of AI agents, thus the need of security and privacy when dealing with tools that, potentially, could harvest personal or company data for a variety of reasons.
More in detail the concept that we were trying to focus on, was the scenario in which an employee, by chance or intentionally, had its AI agent to look at a confidential file which the company doesn't want to be leaked or simply does not trust the AI agent's vendor.

This particular fear, to not trust AI vendors, it's not unfounded; several recent incidents show how sensible data can end up outside the boundaries an organization intended to enforce, or how AI providers themselves may proceed to process user data in ways that raise legitimate concerns.

A first well known example is the 2023 Samsung incident, where some employees pasted proprietary source code into a public LLM chatbot for debugging purposes, effectively transmitting confidential IP (intellectual property) to an external, unintended infrastructure. This incident led Samsung to ban all AI tools altogether\cite{samsung2023leak}

A second, more institutional, example involves OpenAI: in december 2024, the Italian Data Protection Authority fined OpenAI for \texteuro15 million, finding that ChatGPT was processing users' personal data to train underlying models without the legal basis required by the GDPR; not to mention a security breach, that took place in march 2023, which exposed users' data and payment information\cite{garante2024openai}.

Also Anthropic was acting in regard of the guarantees that have to be assumed by a vendor, in particular it had to update its customer terms\cite{anthropic2025terms}. They explicitly stated that each conversation made with Claude could be used for model training, with the exception of Claude for Work, Claude Enterprise, Claude Gov and any API-based access\cite{anthropic2025terms}.

\subsection{Activities done and Results}
Fornire una breve sintesi delle attività che sono state svolte durante il periodo di tesi e dei principali risultati ottenuti (applicazioni o sistemi sviluppati). (circa 1 pagina)

\subsection{Technologies and Tools}
% \section{Tecnologie e strumenti utilizzati per lo sviluppo del progetto}
% Fornire una descrizione delle tecnologie e degli strumenti utilizzati per la realizzazione del progetto (massimo 1 pagina per ogni tecnologia/strumento).